


\section{相关工作}
\label{sec.related_work}


\begin{table}[t]
\centering
\addtolength{\tabcolsep}{-0.3em}
\resizebox{\columnwidth}{!}{
\begin{tabular}{l c c c c c}
\toprule
方法 & 训练 & \makecell{面向目标的\\子智能体} & \makecell{整体式\\编排} & \makecell{MAS\\程度} & \makecell{MAS 与 SAS\\对比分析} \\
\hline
MAS-Zero \cite{Ke2025MASZero} & \xmark & \xmark & \cmark & \xmark & \xmark \\
AFlow \cite{zhang2024aflowautomatingagenticworkflow} & \xmark & \xmark & \xmark & \xmark & \xmark \\
DyLAN \cite{liu2023dynamic} & \xmark & \xmark & \xmark & \xmark & \xmark \\
 \midrule
MAS-GPT \cite{ye2025masgpttrainingllmsbuild} & SFT & \xmark & \cmark & \xmark & \xmark \\
CoRL \cite{jin2025controllingperformancebudgetcentralized} & RL & \xmark & \xmark & \xmark & \xmark \\
Puppet \cite{dang2025multiagentcollaboration} & RL & \xmark & \xmark & \xmark & \xmark \\
xRouter \cite{qian2025xrouter} & RL & \xmark & \xmark & \xmark & \xmark \\
W4S \cite{nie2025weakforstrong} & RL & \xmark & \xmark & \xmark & \xmark \\
ToolOrchestra \cite{su2025toolorchestraelevatingintelligenceefficient} & RL & \cmark & \xmark & \xmark & \xmark \\
\textbf{\ourframework} & RL & \cmark & \cmark & \cmark & \cmark \\
\bottomrule
\end{tabular}
}
\caption{主流自动化 MAS 方法对比。
% \zixuan{or we can have a parato front fig}
}
\label{tab:method_comparison}
\vspace{-2\baselineskip}
\end{table}

\subsection{从单智能体系统到多智能体系统}
尽管\textit{智能体性}通常被视为一个连续谱\cite{kapoor2024ai}，清楚地区分单智能体系统（SAS）与多智能体系统（MAS）仍是一项基础工作。我们将 SAS 定义为仅有一个推理中心且只对应一个目标的系统：所有感知、规划和行动都在由单个 LLM 实例控制的顺序控制循环内执行，以实现该目标。这一定义包括采用工具使用\citep{yao2023react}、自我反思\citep{madaan2024self}或 CoT 推理的系统。相比之下，MAS 由多个具有各自目标的 LLM 智能体组成；它们通过消息传递、共享记忆或显式编排协议等结构化通信机制交互\citep{xi2023risepotentiallargelanguage}。每个智能体都可以维护自己的上下文、目标和工具，系统级行为则通过智能体间的协作以\emph{集体推理}形式涌现，而非源自单一的整体推理过程。

% While MAS is related to other concepts such as modular reasoning or tool composition, its defining characteristic lies in the presence of multiple interacting reasoning entities. \zixuan{there are also some related concepts}

% Early MAS designs were largely manually constructed, with fixed interaction patterns like debate \citep{du2023improvingfactualityreasoning}. More recently, there has been growing interest in \textbf{automatic MAS design}, where the system structure, agent roles, and interactions are generated dynamically. 
根据不同语境下对子智能体的定义，MAS 设计还与 LLM 路由器{\citep{su2025toolorchestraelevatingintelligenceefficient,zhang2025routerr1teachingllmsmultiround}}（将 LLM 视为子智能体）及技能（Skills）\citep{anthropic_agent_skills_2024}（将技能视为子智能体）相关。尽管已有工作研究人工 MAS 设计和自动化推理时编排（见\cref{ap:related}），考虑到\textbf{训练时编排}日益受到关注且与 \ourframework 直接相关，我们将讨论重点放在这一方向。
% \ourframework falls into this latter category, focusing on automatic MAS design that enables scalable and principled orchestration of complex sub-agents.

\subsection{自动化训练时编排}
% \textbf{Inference-time orchestration.} This line of work typically relies on external signals from validation sets or LLM-based judges to guide adaptation \citep{zhou2025multiagentdesignoptimizingagents,zhang2025multiagentarchitecturesearchagentic,hu2025automated,zhang2024aflowautomatingagenticworkflow,Ke2025MASZero}. However, these methods often exhibit limited adaptability due to unreliable adaptation signals.


% these inference-time methods exhibit limited adaptability, as unreliable validation signals or incorrect guidance from LLM-based judges can misdirect the adaptation process.

% MAS-Zero~\citep{Ke2025MASZero} is an exception, performing inference-time MAS design without validation feedback. 

% Recent work on automatic MAS design typically operates at inference time and relies on a validation set to guide adaptation \citep{zhou2025multiagentdesignoptimizingagents,zhang2025multiagentarchitecturesearchagentic,hu2025automated,zhang2024aflowautomatingagenticworkflow}. MAS-Zero~\citep{Ke2025MASZero} differs by performing inference-time MAS design without validation feedback.
% MASS~\citep{zhou2025multiagentdesignoptimizingagents} uses rejection sampling to select effective agent configurations, and MaAS~\citep{zhang2025multiagentarchitecturesearchagentic} extends MASS with a question-wise masking mechanism to adapt subnetworks. ADAS~\citep{hu2025automated} and AFlow~\citep{zhang2024aflowautomatingagenticworkflow} frame MAS design as a code generation task. ADAS stores and searches historical designs based on validation performance, while AFlow enhances this with Monte Carlo Tree Search. 
% One exception is MAS-Zero \citep{Ke2025MASZero}, which performs fully inference-time automatic MAS design without relying on a validation set. 

% \textbf{Training-time orchestration.} 
这一研究方向通过直接训练编排器来避免昂贵的推理时适应。该类别中的多数现有方法采用顺序式编排，并使用多步强化学习进行优化\citep{su2025toolorchestraelevatingintelligenceefficient,nie2025weakforstrong}。
MAS-GPT\citep{ye2025masgpttrainingllmsbuild}是一个例外，它不使用强化学习来训练编排器；然而，其报告性能甚至低于推理时适应方法\citep{Ke2025MASZero}，说明它在学习有效的系统级协作方面存在局限。
据我们所知，\ourframework 是首个执行训练时整体式编排的方法，其中元智能体被训练成在单个决策步骤中生成完整的多智能体系统。这种表述使 \ourframework 的编排器能够在\emph{规划层面而非执行轨迹上}进行推理，从而实现子智能体间的全局协作，避免中间状态引起的误差累积，并让训练目标直接对齐最终任务性能。
% ... + adavtange \zixuan{avoid duplicaate}
表~\ref{tab:method_comparison}系统对比了 \ourframework 与具有代表性的推理时、训练时 MAS 设计方法。结果表明，\ourframework 是唯一同时支持面向目标的子智能体、整体式编排、可配置 \masness，以及旨在理解 MAS 何时及如何优于 SAS 的受控实验的方法。
% \zixuan{need to mention}

% Infernce-time, training time...


